% Source: Weinberg GR, physical PDF pp. 104--110
%         (printed pp. 79--85).
% Coverage: Section 3.5, Time Dilation; equations (3.5.1)--(3.5.6).
% Figures/tables/footnotes: none; reference markers 2--8a and 13.
% Status: source-reviewed and compile-clean.
% Uncertainties: none.

\subsection{Time Dilation}\label{sec:3.5}

Consider a clock in an arbitrary gravitational field, moving with arbitrary
velocity, not necessarily in free fall. The equivalence principle tells us
that its rate is unaffected by the gravitational field if we observe the
clock from a locally inertial coordinate system $\xi^\rho$, so according to
Section~2.2, the spacetime interval $\dd\xi^\rho$ between ticks is governed
in this system by
\[
  \Delta t=
  \left(-\eta_{\rho\sigma}\,\dd\xi^\rho\dd\xi^\sigma\right)^{1/2},
\]
where $\Delta t$ is the period between ticks when the clock is at rest in the
absence of gravitation. Hence in any arbitrary coordinate system the
spacetime interval between ticks will be governed by
\[
  \Delta t=
  \left(
  -\eta_{\rho\sigma}
  \frac{\partial\xi^\rho}{\partial x^\mu}\dd x^\mu
  \frac{\partial\xi^\sigma}{\partial x^\nu}\dd x^\nu
  \right)^{1/2},
\]
or, introducing the metric tensor \eqref{eq:3.2.7},
\[
  \Delta t=\left(-g_{\mu\nu}\,\dd x^\mu\dd x^\nu\right)^{1/2}.
\]
If the clock has velocity $\dd x^\mu/\dd t$, then the time interval $\dd t$
between ticks will be given by
\begin{equation}
  \frac{\dd t}{\Delta t}
  =
  \left(
  -g_{\mu\nu}
  \frac{\dd x^\mu}{\dd t}\frac{\dd x^\nu}{\dd t}
  \right)^{-1/2}.
  \tag{3.5.1}\label{eq:3.5.1}
\end{equation}
In particular, if the clock is at rest this becomes
\begin{equation}
  \frac{\dd t}{\Delta t}=(-g_{00})^{-1/2}.
  \tag{3.5.2}\label{eq:3.5.2}
\end{equation}

We cannot observe the time-dilation factors appearing in \eqref{eq:3.5.1}
and \eqref{eq:3.5.2} by merely measuring the time interval $\dd t$ between
ticks and comparing with the value $\Delta t$ specified by the manufacturer,
because the gravitational field affects our time standards in exactly the
same way as it affects the clock being studied. That is, if our standard
clock says that a certain physical process takes 1 sec at rest in the
absence of gravitation, then it will also tell us that it takes 1 sec in the
presence of gravitation, both standard clock and process being affected by
the field in the same way. However, we can compare the time-dilation factors
at two different points in a field. For instance, suppose that at point 1 we
observe the light coming from a particular atomic transition at point 2. If
points 1 and 2 are at rest in a stationary gravitational field, then the
time taken for a wave crest to travel from 2 to 1 will be a constant, given
by the integral of \eqref{eq:3.2.10} over the path, and therefore the time
between the arrival at point 1 of successive crests will equal the time
$\dd t_2$ between their departure at point 2, given by \eqref{eq:3.5.2} as
\[
  \dd t_2=\Delta t\,[-g_{00}(x_2)]^{-1/2}.
\]
If the same atomic transition occurs at point 1, then the time between crests
of the light waves to be observed at point 1 will be
\[
  \dd t_1=\Delta t\,[-g_{00}(x_1)]^{-1/2}.
\]
Hence, for a given atomic transition, the ratio of the frequency (observed at
point 1) of the light from point 2 to that of the light from point 1 will be
\begin{equation}
  \frac{\nu_2}{\nu_1}
  =\left(\frac{g_{00}(x_2)}{g_{00}(x_1)}\right)^{1/2}.
  \tag{3.5.3}\label{eq:3.5.3}
\end{equation}
In the weak-field limit $g_{00}\simeq-1-2\phi$ and $\phi\ll1$, so
$\nu_2/\nu_1=1+\Delta\nu/\nu$, where
\begin{equation}
  \frac{\Delta\nu}{\nu}=\phi(x_2)-\phi(x_1).
  \tag{3.5.4}\label{eq:3.5.4}
\end{equation}
(For a \emph{uniform} gravitational field, this result could be derived
directly from the Principle of Equivalence without introducing a metric or
affine connection.)

Let us apply Eq.~\eqref{eq:3.5.4} to the case of light from the sun's surface
observed on the earth. The sun's gravitational potential can be calculated
as
\[
  \phi_\odot=-\frac{GM_\odot}{R_\odot},
\]
where $M_\odot$ and $R_\odot$ are the sun's mass and radius,
\[
  M_\odot=1.97\times10^{33}\ \mathrm{g},
  \qquad
  R_\odot=0.695\times10^6\ \mathrm{km},
\]
and $G$ is the gravitational constant
\begin{equation}
  G=6.67\times10^{-8}\ \frac{\mathrm{erg\,cm}}{\mathrm{g}^2}
  =7.41\times10^{-29}\ \frac{\mathrm{cm}}{\mathrm{g}}.
  \tag{3.5.5}\label{eq:3.5.5}
\end{equation}
(Here we have used our convention that $c=1$ to set
$1\ \mathrm{sec}=3\times10^{10}\ \mathrm{cm}$; in c.g.s.\ units the
quantity $7.41\times10^{-29}\ \mathrm{cm/g}$ would have to be called
$G/c^2$.) We find that the potential on the surface of the sun is
\[
  \phi_\odot=-2.12\times10^{-6}.
\]
The gravitational potential of the earth is negligible in comparison, so
ideally the frequency of light from the sun should be shifted to the red by
2.12 parts per million as compared with light emitted by terrestrial atoms.

The difficulty in measuring the solar gravitational red shift can be
appreciated if we reflect that a motion of the source by a velocity $v$ along
the earth--sun direction will produce an additional Doppler frequency shift
$\Delta\nu/\nu=v$ [recall Eq.~(2.2.2)], so the gravitational red shift can
be masked by a velocity $v=2\times10^{-6}$, or in c.g.s.\ units,
$v=0.6\ \mathrm{km/sec}$. It is not the rotation of the earth or the sun
that bothers us; these are known effects which can easily be taken into
account. Thermal effects are more serious; at a temperature of
$3000^\circ\mathrm K$ the thermal velocity of typical light elements
($\mathrm C,\mathrm N,\mathrm O$) is about $2\ \mathrm{km/sec}$, giving a
Doppler broadening about three times larger than the expected gravitational
red shift. However, thermal motions only broaden lines, not shift them, so
they too can be lived with. The really bothersome problems arise from unknown
Doppler shifts owing to the convection of gases in the solar atmosphere. In
fact, the frequency shift is observed to vary from place to place on the
solar disk, and is occasionally even toward the blue! The convection tends
to be vertical, so we can minimize the Doppler shifts it produces by looking
at the limb of the sun, where the motion is mostly at right angles to the
line of sight. Until recently the best result that could be achieved in this
way was that the solar gravitational red shift is of the order of 2 parts
per million.~\hyperref[ch3-ref-2]{[2]} In the last few years improved
observational techniques~\hyperref[ch3-ref-3]{[3]} have yielded a much
better value of the red shift, equal to $1.05\pm0.05$ times the predicted
value. However, it is too early to say that this result closes the story, at
least until it can be corroborated.

The red shifts are much larger for white dwarf stars like Sirius B and 40
Eridani B. Such stars have masses typically of the order of one solar mass,
and radii of the order of $1/10$ to $1/100$ the sun's radius, so the red
shift of spectral lines from their surface is about 10 to 100 times greater
than for the sun, or roughly 1 part in $10^4$ to $10^5$. Although this
alleviates problems arising from convective Doppler shifts or temperature or
pressure, a new difficulty enters here: It is difficult to determine the
value of the gravitational potential $\phi$ with which to compare the
measured value of $\Delta\nu/\nu$. If we know the mass of a white dwarf
star, we can deduce a rough value for its radius and surface gravitational
potential from astrophysical theory,~\hyperref[ch3-ref-4]{[4]} but the only
white dwarf stars whose mass can be measured are members of binary star
pairs. For instance the mass of Sirius B is determined by calculating the
total mass of Sirius A and B from their separation and period, and then
subtracting the mass of Sirius A calculated from stellar theory. However,
the scattering of light from Sirius A by the atmosphere of Sirius B makes
the gravitational red shift on Sirius B very difficult to
measure.~\hyperref[ch3-ref-13]{[13]} On the other hand, 40 Eridani B is
sufficiently far from 40 Eridani A so that scattering of light is no
problem, and the mass of 40 Eridani B can be determined separately from that
of A by locating their center of mass in addition to measuring their period
and separation. However, because 40 Eridani B and A are so far apart, their
period of revolution is very long, and there has not yet been time to
determine the mass of B very accurately. The best predicted value of the
surface gravitational potential is
\[
  \phi=-(5.7\pm1)\times10^{-5},
\]
in good agreement with the observed~\hyperref[ch3-ref-5]{[5]} red shift
\[
  \frac{\Delta\nu}{\nu}=-(7\pm1)\times10^{-5}.
\]
Taking account of Stark shifts in the spectrum of 40 Eridani B appears to
improve the agreement.~\hyperref[ch3-ref-5a]{[5a]}

The empirical evidence for the red shift predicted by the Principle of
Equivalence was much improved in 1960, by a terrestrial experiment performed
by Pound and Rebka.~\hyperref[ch3-ref-6]{[6]} They allowed a $\gamma$ ray
emitted by a $14.4\ \mathrm{keV}$, $0.1\ \mu\mathrm{sec}$ transition in
$\mathrm{Fe}^{57}$ to fall $22.6\ \mathrm m$, and observed its resonant
absorption by an $\mathrm{Fe}^{57}$ target. (Normally \emph{resonant}
absorption is impossible for such a narrow $\gamma$-ray line, because the
recoil of the emitting nucleus lowers the $\gamma$-ray energy below the
nuclear energy difference, whereas to produce the inverse transition in the
target nucleus, which also recoils, a little more energy than the nuclear
energy difference is needed. This experiment was made possible by the
M\"ossbauer effect,~\hyperref[ch3-ref-7]{[7]} in which the recoil momentum on
emission and absorption is taken up by the whole crystal, so that essentially
no energy is lost to recoil on emission or absorption.) The difference in
the gravitational potential from top to bottom is
\begin{align*}
  \Delta\phi
  =\phi_{\mathrm{top}}-\phi_{\mathrm{bottom}}
  &=-\frac{(980\ \mathrm{cm/sec^2})(2260\ \mathrm{cm})}
  {(3\times10^{10}\ \mathrm{cm/sec})^2}
  \\
  &=-2.46\times10^{-15},
\end{align*}
and if the equivalence principle is correct we would expect the photon
arriving at the target to be shifted upward in frequency by an amount
$\Delta\nu/\nu=-\Delta\phi$, lowering the counting rate by a factor
\[
  C=\frac{\Gamma^2}{\Delta\nu^2+\Gamma^2},
\]
where $\Gamma$ is the full width of the $\gamma$-ray line at half maximum.
(Note that $\Gamma$ appears here rather than $\Gamma/2$, because we have to
fold together an emission coefficient proportional to
$[(\nu+\Delta\nu)^2+(\Gamma/2)^2]^{-1}$ with an absorption coefficient
proportional to $[\nu^2+(\Gamma/2)^2]^{-1}$.) But in this transition the
fractional width was $\Gamma/\nu=1.13\times10^{-12}$, which is larger than
the predicted $\Delta\nu/\nu$ by a factor of 460, so the reduction in the
counting rate was by only 1 part in $2.1\times10^5$! This would seem to make
the experiment impossible, and indeed Pound and Rebka had originally
thought that they might have to let the $\gamma$ ray fall several kilometers
in order to get a frequency shift $\Delta\nu$ comparable with $\Gamma$, but
happily they thought of a trick which let them measure very small frequency
shifts. Their idea was to move the $\gamma$-ray source up and down with
velocity $v_0\cos\omega t$, where $\omega$ is some arbitrary fixed frequency
(10--50 cps) and $v_0$ is also arbitrary, but much greater than
$-\Delta\phi$, that is, much greater than
$7.4\times10^{-5}\ \mathrm{cm/sec}$. To the gravitational violet shift
$\Delta\nu_G$ there is then added a larger Doppler shift
$\Delta\nu_D/\nu=-v_0\cos\omega t$ (see Section~2.2), so the counting rate
is reduced by a time-dependent factor
\begin{align*}
  C(t)
  &=\frac{\Gamma^2}{(\Delta\nu_G+\Delta\nu_D)^2+\Gamma^2}
  =
  \frac{(\Gamma/\nu)^2}
  {(\Delta\nu_G/\nu-v_0\cos\omega t)^2+(\Gamma/\nu)^2}
  \\
  &\simeq
  \frac{(\Gamma/\nu)^2}
  {v_0^2\cos^2\omega t+(\Gamma/\nu)^2}
  \left\{1+
  \frac{2(\Delta\nu_G/\nu)v_0\cos\omega t}
  {v_0^2\cos^2\omega t+(\Gamma/\nu)^2}\right\},
\end{align*}
and $\Delta\nu_G$ can be picked out by looking for a term linear in
$\cos\omega t$, for instance, by measuring the asymmetry between the number
of counts registered when the source is going up (for example,
$\cos\omega t>1/\sqrt2$) and down ($\cos\omega t<-1/\sqrt2$). In this way
Pound and Rebka obtained a value for $\Delta\nu_G/\nu$ about four times
larger than the expected value $2.46\times10^{-15}$. This discrepancy was
actually an intrinsic frequency shift owing to the difference between the
source and target crystals (including differences in their temperature) and
was removed by subtracting the asymmetry in $\gamma$-ray counts when the
source is below the target from the asymmetry when the target is below the
source. The final result for the gravitational frequency shift was
$\Delta\nu/\nu=(2.57\pm0.26)\times10^{-15}$, in excellent agreement with
the predicted value $2.46\times10^{-15}$. The agreement between theory and
experiment has since~\hyperref[ch3-ref-8]{[8]} been improved to about 1
percent.

There have also been proposals~\hyperref[ch3-ref-8a]{[8a]} to measure the
gravitational red shift of light from an artificial satellite. At a point
directly below perigee there is no first-order Doppler shift because the time
for the light to reach us from the satellite is momentarily constant. In
this case the frequency shift of the emitted light must be determined from
\eqref{eq:3.5.1}, whereas the frequency shift of our laboratory time
standards may be calculated from \eqref{eq:3.5.2}, if we ignore the rotation
of the earth. It follows that the frequency $\nu_s$ of a given atomic line
from the satellite will be related to the frequency $\nu_e$ of the same line
on earth by
\begin{equation}
  \frac{\nu_s}{\nu_e}
  =
  \frac{\left(
  -g_{\mu\nu}\dfrac{\dd x^\mu}{\dd t}
  \dfrac{\dd x^\nu}{\dd t}\right)_s^{1/2}}
  {(-g_{00})_\oplus^{1/2}}.
  \tag{3.5.6}\label{eq:3.5.6}
\end{equation}
The velocity $v_s$ of the satellite is given by
\[
  v_s^2=-\phi_s=\frac{GM_\oplus}{R_\oplus+H},
\]
where $H$ is the altitude of the satellite and $M_\oplus$ and $R_\oplus$
are the mass and radius of the earth,
\[
  M_\oplus=5.983\times10^{27}\ \mathrm g,
  \qquad
  R_\oplus=6.371\times10^8\ \mathrm{cm}.
\]
In the weak-field approximation we have
\begin{align*}
  \left(
  -g_{\mu\nu}\frac{\dd x^\mu}{\dd t}
  \frac{\dd x^\nu}{\dd t}\right)_s
  &\simeq -(g_{00})_s-v_s^2
  =1+2\phi_s-v_s^2
  \\
  &\simeq 1-\frac{3GM_\oplus}{R_\oplus+H},
\end{align*}
and
\[
  (-g_{00})_\oplus
  \simeq1+2\phi_\oplus
  \simeq1-\frac{2GM_\oplus}{R_\oplus}.
\]
So to this order Eq.~\eqref{eq:3.5.6} gives a frequency ratio
$\nu_s/\nu_e=1+\Delta\nu/\nu$, where
\begin{align*}
  \frac{\Delta\nu}{\nu}
  &=-\frac32\frac{GM_\oplus}{R_\oplus+H}
  +\frac{GM_\oplus}{R_\oplus}
  \\
  &\simeq-3.47\times10^{-10}
  \left\{\frac{3R_\oplus}{R_\oplus+H}-2\right\}.
\end{align*}
We see that at low altitudes there is a purely special-relativistic red shift
(see Section~2.2), to which is added at higher altitudes a
general-relativistic violet shift, yielding a net red shift for
$H<R_\oplus/2$ and a net violet shift for $H>R_\oplus/2$.

Incidentally, the gravitational red shift of light rising from a lower to a
higher gravitational potential can to some extent be understood as a
consequence of quantum theory, energy conservation, and the ``weak''
Principle of Equivalence. When a photon is produced at point 1 by some heavy
nonrelativistic apparatus, an observer in a locally inertial coordinate
system moving with the apparatus will see its internal energy and hence its
inertial mass change by an amount related to the photon frequency $\nu_1$ he
observes, that is, by
\[
  \Delta m_1=-h\nu_1,
\]
where $h=6.625\times10^{-27}\ \mathrm{erg\,sec}$ is Planck's constant.
Suppose that the photon is then absorbed at point 2 by a second heavy
apparatus; an observer in a freely falling system will see the apparatus
change in inertial mass by an amount related to the photon frequency $\nu_2$
he observes, that is, by
\[
  \Delta m_2=h\nu_2.
\]
However, the total internal plus gravitational potential energy of the two
pieces of apparatus must be the same before and after these events, so
\[
  0=\Delta m_1+\phi_1\Delta m_1+\Delta m_2+\phi_2\Delta m_2,
\]
and therefore
\[
  \frac{\nu_2}{\nu_1}
  =\frac{1+\phi_1}{1+\phi_2}
  \simeq1+\phi_1-\phi_2,
\]
in agreement with our previous result. (Also, it makes no difference whether
the photon frequencies are measured in locally inertial systems, because
the gravitational field in any other frame will affect the rate of the
observer's standard clock in the same way as it affects the $\nu$'s.) This
result can be interpreted as saying that a photon in a gravitational field
has ``kinetic energy'' $h\nu$ and ``potential energy'' $h\nu\phi$, their sum
remaining constant. However, I have insisted on including a nonrelativistic
emitter and absorber in the above calculation, because the concept of
gravitational potential energy for a photon is otherwise without foundation.

This derivation rests on the Principle of Equivalence in \emph{two}
respects: It assumes that the change in gravitational mass of the apparatus
equals the change in its inertial mass and hence its internal energy; and it
also assumes that in a freely falling frame the relation between photon
energy and frequency is unaffected by the presence of gravitational fields.
Hence even if we suppose that the E\"otv\"os--Dicke experiments could improve
to an unlimited accuracy, and that gravitational mass were found to equal
inertial mass exactly, still there would be some point in verifying the
gravitational red shift of spectral lines, as an independent test of the
Principle of Equivalence.
